% Include only after the full GPU benchmark passes its completeness audit.
\section{Expanded Offline Benchmark}
\label{sec:expanded}
\subsection{Question and Fixed Protocol}
The preceding streaming failure changes both the learning regime and the
selection algorithm. We therefore evaluate the original offline scoring rule
on a broader, separately specified benchmark. This study asks whether global
top-$K$ selection by the sum of standardized entropy and squared gradient norm
improves final competence beyond random, entropy-only, and representative
storage. It does not modify the gate using the new test outcomes.

The locally recorded protocol and its source hashes precede the full run; it
is not an externally registered preregistration. Ten paired seeds (100--109)
randomize initialization, new-example order, random selection, and the partition
of ten classes into two successive five-class tasks. Digits uses a stratified
80/20 split per seed; Fashion-MNIST and CIFAR-10 use their full official training
and test splits. CIFAR-10~\cite{cifar} contains natural color images, broadening the domains
beyond grayscale digits and clothing. Its Kaggle-hosted mirror is validated
against torchvision's official batch-file checksums. No leaderboard score is
used as a comparator.

All selectors within a seed share the Task-A checkpoint. They store examples
from its complete training pool, then train on Task B with a fixed replay
exposure. Task identity is never supplied at inference: the classifier always
predicts among all ten classes. New-example shuffling and replay-position draws
use separate random generators, unlike the coupled generator in the archived
implementation. Thus even the no-replay anchor receives the same new-example
order. These changes make the expansion a generalization study, not an exact
reproduction of the original data split and training trajectory.

Digits retains the original 64-hidden-unit MLP, SGD learning rate 0.5, 60/40
Task-A/Task-B epochs, and batches containing 32 new and two replay examples.
Its storage budgets are 20, 40, and 100. Fashion-MNIST uses a 128-hidden-unit
MLP, learning rate 0.1, 10/10 epochs, and 128 new plus eight replay examples.
Its budgets are 50, 200, and 500. CIFAR-10 uses three convolution--ReLU--pooling
blocks with 32, 64, and 128 channels, followed by a 128-unit hidden layer and
ten outputs. Each block uses a $3\times3$ convolution with unit padding and
$2\times2$ max pooling. It uses SGD at 0.05 with momentum 0.9, 20/20 epochs, the same
128:8 batch composition, and budgets 50, 200, and 500. Pixels are normalized
by channel; there is no augmentation, batch normalization, pretrained backbone,
learning-rate search, or test-based checkpoint selection. Accordingly, the
CNN is a controlled learner rather than a competitive image-classification
submission.

\subsection{Selectors and Resource Accounting}
Nine bounded selectors are evaluated: random, class-balanced random, Doubt,
Curiosity, EgoGate, embedding k-center, the existing EgoGate-diverse rule,
highest training loss, and a class-balanced herding adaptation. The latter
greedily approximates each class's mean using normalized hidden features and
unique exemplars. It adapts the exemplar-selection component associated with
iCaRL~\cite{icarl}; it does not implement iCaRL's distillation or nearest-exemplar-mean
classifier. No-replay and full-pool replay provide separately identified anchors,
not equal-memory competitors. Full-pool replay uses the same replay count per
update and is not assumed to be an accuracy upper bound.

The MLP gradient norm is computed analytically. The CNN score sums the squared
components of exact per-example parameter gradients, evaluated with vectorized
automatic differentiation in batches of 16. Both use training labels and
global pool standardization, preserving the original offline score rather than
the streaming bounded transform. All selection uses the pre-Task-B model.

The bounded methods have identical optimizer-update and training-exposure
counts within each dataset/seed group. Selection and Task-B training times are
measured separately with CUDA synchronization. Persistent exemplar bytes,
selected-class counts, selected-index hashes, and peak CUDA allocations are
recorded. Peak allocations include the resident datasets and models; they are
not isolated incremental buffer cost. CPU selection scratch space is not
included in CUDA allocations. The full candidate pool is available during
selection, so this experiment does not claim bounded online working memory.
Methods execute in randomized order within each seed. Independent seeds may
run on separate GPUs; all paired methods within a seed use the same device.

\subsection{Endpoints and Statistical Analysis}
The primary endpoint is final mean accuracy, weighting the old and new task
accuracies equally. This prevents apparent retention improvements from being
interpreted without measuring new-task learning. The 27 specified primary
contrasts compare EgoGate with random, Doubt, and herding in each of the nine
dataset/budget cells. Two-sided paired $t$ tests use one Holm-corrected family.
We report pointwise 95\% intervals for seed-paired differences; these intervals
are not simultaneous confidence intervals. Forgetting, task-specific accuracy,
other selector comparisons, coverage, and costs are secondary outcomes.
The seeds vary several nuisance factors, not independent datasets. There is
no significance-based stopping, winner-only reporting, or inference from the
reduced execution-check run. Complete-grid and exposure audits precede analysis.
